\subsection{Vergleich mit alternativen Template-Strategien}
\label{subsec:vergleich-template-strategien}

Verglichen werden Architekturstrategien anhand von Wiederverwendung,
Konsistenz, Variabilität, Einstiegskomplexität, Pflege, Testbarkeit und
Updateübernahme
\parencite{krueger1992software-reuse,clements2001software-product-lines}.
Die qualitativen Stufen bilden kein Gesamtranking; hohe Einstiegskomplexität
ist beispielsweise nachteilig.

\emph{Separate Template-Repositories} bleiben einzeln übersichtlich, können
aber divergieren. Ein \emph{Copy-Paste-Starter} minimiert die Vorlagenlogik,
liefert jedoch keinen Rückkanal für spätere Verbesserungen. GitHub-Templates
erzeugen entsprechend eine eigene, nicht verwandte Historie
\parencite{github2026repository-template}.
Separate Repositories gewichten damit organisatorische Isolation höher als die
zentrale Konsistenz gemeinsamer Bausteine.

\emph{Framework-spezifische Starter} wie \texttt{create-vite} bieten einen
einfachen Einstieg in einen engen Ausschnitt. Backend, Desktop, Persistenz und
Gesamtworkflow müssen bei Bedarf separat ergänzt werden
\parencite{vite2026guide}.
Für ein kleines homogenes Frontend kann diese geringere Einstiegskomplexität
angemessener sein als der untersuchte Full-Stack-Ansatz.

Ein \emph{monolithisches Universal-Template} vereinfacht die Auswahl, belastet
kleine Projekte aber mit ungenutzten Komponenten. Das untersuchte
\emph{Single-Repository mit Profilen und Capabilities} kombiniert zentrale
Pflege mit ausgewählten Pfaden. Dafür steigen Master-, Resolver- und
Matrixkomplexität; Updates erreichen generierte Projekte nicht automatisch.
Es verbindet somit zentrale Variabilität vor der Generierung mit der
Eigenständigkeit eines Copy-Starters danach.

\begin{table}[htbp]
  \centering
  \scriptsize
  \setlength{\tabcolsep}{3pt}
  \renewcommand{\arraystretch}{1.08}
  \caption{Qualitativer Vergleich von Template-Strategien}
  \label{tab:template-strategy-comparison}
  \begin{tabularx}{\textwidth}{@{}>{\raggedright\arraybackslash}p{0.18\textwidth}>{\centering\arraybackslash}p{0.12\textwidth}>{\centering\arraybackslash}p{0.11\textwidth}>{\centering\arraybackslash}p{0.12\textwidth}>{\centering\arraybackslash}p{0.13\textwidth}>{\centering\arraybackslash}p{0.12\textwidth}>{\raggedright\arraybackslash}X@{}}
    \toprule
    \textbf{Strategie} & \textbf{fortlaufende Wiederverwendung} &
      \textbf{Konsistenz} & \textbf{Variabilität} &
      \textbf{Einstiegs\-komplexität} & \textbf{Pflege\-bündelung} &
      \textbf{Updates bestehender Projekte} \\
    \midrule
    Separate Template-Repositories & mittel & mittel & mittel & niedrig &
      niedrig & je Vorlage beziehungsweise Projekt manuell \\
    Copy-Paste-Starter & niedrig nach Kopie & niedrig & hoch & niedrig &
      niedrig & keine automatische Übernahme \\
    Framework-spezifischer Starter & hoch im engen Scope & hoch im engen Scope
      & niedrig bis mittel & niedrig & hoch im Starter & meist manuell nach
      Generierung \\
    Monolithisches Universal-Template & mittel & hoch & niedrig & hoch & hoch
      & abhängig vom gewählten Ableitungsweg \\
    Single-Repository mit Profilen & hoch im definierten Scope & hoch & hoch,
      kontrolliert & mittel & hoch & im Master zentral, in abgeleiteten
      Projekten manuell \\
    \bottomrule
  \end{tabularx}
  \smallskip
  {\footnotesize Quelle: Eigene qualitative Gegenüberstellung auf Grundlage
  von \textcite{krueger1992software-reuse},
  \textcite{clements2001software-product-lines},
  \textcite{github2026repository-template} und \textcite{vite2026guide}.\par}
\end{table}


Der profilbasierte Ansatz passt vor allem zu wiederkehrenden, technisch
verwandten Vorhaben. Kleine Einzelprojekte können von Framework-Startern,
organisatorisch getrennte Produktlinien von separaten Repositories profitieren.
Es gibt keinen universellen Sieger, sondern unterschiedliche Trade-offs.
\beginnerinput{workspace/statement/700_bewertung/750}
